\section{Introdução}

Em um aplicativo de transporte urbano, cada solicitação precisa ser associada a um veículo livre. A distância que esse veículo percorre até a origem do passageiro (aproximação, ou \textit{pickup}) é quilometragem sem passageiro, aqui chamada \textit{km vazio}. O km vazio consome tempo e combustível sem gerar corrida e, por isso, é um alvo natural de otimização. A distância de aproximação depende de dois fatos: onde cada veículo está e qual veículo é escolhido para cada pedido. Quando vários pedidos chegam em curto intervalo e os veículos livres são poucos, escolher o mais próximo para cada pedido, um de cada vez, pode consumir justamente o veículo de que outro pedido precisaria mais.

Este trabalho investiga uma alternativa localizada nessa decisão: reunir os pedidos já recebidos em lotes de 60 segundos e atribuí-los, em conjunto, aos veículos elegíveis de modo que o maior número possível de pedidos seja atendido e a soma das distâncias de aproximação seja mínima. Essa atribuição conjunta é resolvida pelo algoritmo Húngaro \cite{kuhn1955, munkres1957}. O foco é a atribuição de pedidos já recebidos; não é previsão de demanda, não é reposicionamento preventivo e não é um plano diário global.

A pergunta de pesquisa é: \textbf{como a atribuição conjunta de solicitações em lotes de 60 segundos, pelo algoritmo Húngaro, afeta os deslocamentos sem passageiros e o atendimento em comparação com políticas de despacho pelo veículo mais próximo?} O objetivo geral é medir esse efeito em duas linhas de avaliação separadas e nunca somadas: (A) uma \textbf{avaliação contrafactual por simulação com entradas históricas} de uma operação real da Grande Vitória (ES), e (B) uma \textbf{avaliação em cenários sintéticos} construídos sobre a mesma região, de caráter exploratório e complementar. Os objetivos específicos são: (i) formalizar a atribuição por lote; (ii) implementar o método e verificar sua otimalidade por lote contra um oráculo independente; (iii) comparar, com os mesmos pedidos, frota, jornadas e distâncias, o Húngaro com o despacho pelo mais próximo imediato e em lotes; (iv) avaliar km vazio, atendimento e espera; e (v) delimitar as condições e as limitações dos resultados.

As hipóteses, fixadas antes da execução nos períodos reservados, são duas e são avaliadas separadamente. \textbf{H1 (redução):} o km vazio com atribuição conjunta é menor que o do despacho imediato (contraste principal) e que o do despacho em lotes pelo mais próximo (complementar). \textbf{H2 (preservação do atendimento):} a redução só é considerada com serviço preservado se não houver violação de compromisso causada pela política, o número de atendidos não for menor que o da referência no período e o percentil 90 (P90) da espera não piorar. A hipótese de redução é uma hipótese, não uma promessa; atendimento e cobertura funcionam como controles da comparação.

Um estudo anterior dos autores avaliou o reposicionamento preventivo de veículos na véspera (horizonte D+1, com algoritmo guloso). Esses métodos não reduziram o km vazio nos períodos avaliados, e o resultado negativo permanece válido; ele responde a outra pergunta (posicionar antes) e não é objeto deste artigo, exceto como antecedente \cite{pavone2012, kontou2020, pouls2025}.

A contribuição é empírica e metodológica: uma formulação explícita da atribuição por lote, uma implementação verificada, uma avaliação com configuração congelada antes de dois períodos nunca simulados e o relato de um resultado desfavorável (o atendimento em outubro) junto ao favorável. Os resultados valem para o simulador, os dados e a região descritos e não equivalem à economia na operação real.

A Seção 2 apresenta a fundamentação e o algoritmo; a Seção 3, os dados, as políticas, o desenho e as métricas; a Seção 4, os resultados e a discussão; a Seção 5, as conclusões.
